\chapter{مرجع کامل توابع درایورها}

\section{ذخیره شناسه حسگر}

\subsection{تابع \lr{crc8\_maxim}}
ورودی آن اشاره‌گر داده و طول است و
\lr{CRC}
هشت‌بیتی بازتاب‌یافته با چندجمله‌ای
\lr{0x8C}
را برمی‌گرداند. تابع فقط داخلی است و برای هفت بایت نخست شناسه یک‌سیمه استفاده می‌شود.

\subsection{تابع \lr{crc32\_record}}
روی تمام بایت‌های رکورد به‌جز فیلد
\lr{CRC}
نهایی اجرا می‌شود. مقدار آغازین همه یک، چندجمله‌ای بازتاب‌یافته
\lr{0xEDB88320}
و خروجی معکوس‌شده است. هدف تشخیص رکورد ناقص، فرسوده یا متعلق به نسخه دیگر است.

\subsection{تابع \lr{SensorAddressStore\_IsValidRom}}
یک شناسه هشت‌بایتی می‌گیرد. اشاره‌گر تهی، کد خانواده غیر از ۲۸ یا
\lr{CRC}
نامنطبق نتیجه نادرست می‌دهد. این تابع هم برای فرمان جدید رایانه و هم اعتبارسنجی کل جدول استفاده می‌شود.

\subsection{تابع \lr{table\_is\_valid}}
تمام ۳۲ ردیف جدول را به تابع اعتبارسنجی شناسه می‌دهد و با نخستین خطا بازمی‌گردد. سیاست فعلی جدول ناقص را نمی‌پذیرد؛ بنابراین تعویض یک حسگر تنها وقتی ذخیره می‌شود که همه ردیف‌های دیگر نیز معتبر باشند.

\subsection{تابع \lr{SensorAddressStore\_Load}}
رکورد صفحه پایانی
\lr{Flash}
را بدون کپی اولیه بررسی می‌کند. امضا، نسخه، تعداد،
\lr{CRC}
رکورد و
\lr{CRC}
تک‌تک شناسه‌ها باید درست باشند. در موفقیت، جدول به خروجی کپی می‌شود. در شکست، جدول خام قدیمی بررسی می‌شود و در صورت اعتبار، علاوه بر کپی خروجی به قالب جدید ذخیره می‌شود. اگر هر دو نامعتبر باشند نتیجه نادرست است و لایه کاربرد از پیش‌فرض‌ها استفاده می‌کند.

\subsection{تابع \lr{SensorAddressStore\_Save}}
ابتدا کل جدول ورودی را اعتبارسنجی و یک رکورد محلی می‌سازد. سپس
\lr{Flash}
را باز، دقیقاً یک صفحه را پاک و رکورد را کلمه‌به‌کلمه می‌نویسد و در پایان
\lr{Flash}
را قفل می‌کند. هر خطای
\lr{HAL}
نتیجه نادرست می‌دهد. چون پاک‌کردن صفحه برگشت‌ناپذیر و نوشتن مسدودکننده است، تابع فقط در حلقه اصلی فراخوانی می‌شود.

\section{درایور \lr{ADE7880}}

\warningbox{این بخش رفتار نسخه قدیمی درایور را برای سابقه بازآرایی ثبت می‌کند. مرجع \lr{API} فعال و اصلاح‌شده، فصل ۱۵ و فایل \filepath{Libraries/ADE7880/Inc/ade7880.h} است.}

\subsection{تابع \lr{ADE7880\_ReadRegister}}
ورودی‌ها نمونه درایور، نشانی ۱۶بیتی رجیستر، تعداد بایت و مهلت
\lr{HAL}
هستند. بافر چهار‌بایتی داخلی صفر، فرمان خواندن و نشانی ارسال، سپس یک تا چهار بایت دریافت می‌شود. انتخاب تراشه پیرامون کل تراکنش پایین و بالا می‌رود. داده در
\lr{TEMP}
نمونه درایور قرار می‌گیرد و وضعیت خطای
\lr{HAL}
هنوز به فراخواننده بازگردانده نمی‌شود.

\subsection{تابع \lr{ADE7880\_WriteRegister}}
داده ۳۲بیتی را متناسب با طول یک، دو یا چهار بایت با ترتیب پرارزش‌ترین بایت آماده می‌کند، فرمان نوشتن و نشانی را پیشوند می‌دهد و تراکنش
\lr{SPI}
را زیر انتخاب تراشه انجام می‌دهد. طول‌های دیگر نباید ارسال شوند و در نسخه آینده بهتر است نوع نتیجه برای خطا افزوده شود.

\subsection{تابع \lr{ADE7880\_ConfigSPI}}
دسته
\lr{SPI}
و پایه انتخاب تراشه را در نمونه درایور ذخیره می‌کند. سه بار انتخاب تراشه را تغییر می‌دهد تا رابط سریال مدار روی
\lr{SPI}
قرار گیرد، سپس رجیستر پیکربندی دوم را می‌نویسد تا انتخاب قفل شود. پس از بازنشانی سخت‌افزاری باید دوباره اجرا شود.

\subsection{تابع \lr{ADE7880\_Run}}
رجیستر اجرای پردازشگر داخلی مدار را با مقدار یک می‌نویسد. ورودی فقط نمونه درایور است و پس از ارسال یک میلی‌ثانیه صبر می‌کند.

\subsection{تابع \lr{ADE7880\_ReadVRMS}}
سه رجیستر ولتاژ مؤثر فازهای
\lr{A}
تا
\lr{C}
را به‌ترتیب می‌خواند. هر مقدار ۲۴بیتی از سه بایت پایانی بافر ساخته و در ساختار فاز ذخیره می‌شود. تبدیل به ولت در این تابع انجام نمی‌شود.

\subsection{تابع \lr{ADE7880\_ReadIRMS}}
مانند تابع ولتاژ است، با این تفاوت که سه جریان فاز و جریان نول را تازه می‌کند. خروجی‌ها هنوز شمار خام رجیستر هستند.

\subsection{تابع \lr{ADE7880\_ReadPowerz}}
توان حقیقی و توان ظاهری هر سه فاز را از شش رجیستر ۳۲بیتی می‌خواند. داده در اعضای
\lr{WATT}
و
\lr{VA}
قرار می‌گیرد. علامت توان حقیقی باید مطابق قالب دیتاشیت تفسیر شود؛ ساختار از نوع علامت‌دار استفاده می‌کند.

\subsection{تابع \lr{ADE7880\_ReadFPower}}
چهار رجیستر توان حقیقی، راکتیو، ظاهری و ضریب توان بنیادی را می‌خواند و براساس ورودی انتخاب فاز در ساختار
\lr{A}،
\lr{B}
یا
\lr{C}
می‌گذارد. این رجیسترها اشتراکی‌اند؛ بنابراین فراخواننده باید پیش‌تر بیت انتخاب فاز رجیستر پیکربندی را تنظیم کرده باشد. برنامه فعلی این تابع را استفاده نمی‌کند.

\subsection{تابع \lr{ADE7880\_Reset}}
بیت بازنشانی نرم‌افزاری را در رجیستر پیکربندی می‌نویسد و همان رجیستر را تا صفرشدن بیت بررسی می‌کند. حلقه مهلت کلی ندارد؛ قطع ارتباط
\lr{SPI}
می‌تواند راه‌اندازی را متوقف کند. افزودن شمارنده تلاش یکی از اصلاحات ضروری آینده است.

\section{درایور دوگانه \lr{DS18B20}}

\subsection{تابع \lr{DS18B20\_Init}}
هر دو پایه را خروجی تخلیه‌باز با حالت آزاد بالا می‌کند. این تابع باید پیش از نخستین پالس بازنشانی اجرا شود تا پایه پوش‌پول خط یک‌سیمه را در سطح بالا درایو نکند.

\subsection{توابع زمان و شیار}
تابع
\lr{delayus}
با یک حلقه وابسته به کلاک تقریباً به‌اندازه ورودی میکروثانیه مکث می‌کند. توابع
\lr{LH\_signal}
و
\lr{LH\_signal\_1}
بخش پایین و آزاد هر شیار را روی گذرگاه متناظر تولید می‌کنند. این روش در برابر تغییر سطح بهینه‌سازی یا کلاک حساس است.

\subsection{توابع نوشتن}
توابع
\lr{write\_bit}
و
\lr{write\_bit\_1}
براساس صفر یا یک، طول بخش پایین شیار را انتخاب می‌کنند. توابع
\lr{write\_byte}
و
\lr{write\_byte\_1}
هشت بیت را از کم‌ارزش‌ترین بیت می‌فرستند. ورودی بایت پس از ارسال تغییر نمی‌کند.

\subsection{توابع تغییر جهت پایه}
توابع داخلی
\lr{B9\_as\_INPUT}
و
\lr{C7\_as\_INPUT}
خط متناظر را برای نمونه‌برداری ورودی می‌کنند. توابع
\lr{B9\_as\_OUTPUT}
و
\lr{C7\_as\_OUTPUT}
آن را به خروجی تخلیه‌باز بازمی‌گردانند. نام نخست تاریخی است و با پایه واقعی
\lr{PC6}
هم‌خوان نیست؛ تغییر نام باید در یک اصلاح مستقل انجام شود.

\subsection{توابع خواندن}
توابع
\lr{read\_bit}
و
\lr{read\_bit\_1}
شیار خواندن ایجاد، پایه را ورودی، سطح را نمونه و سپس خروجی می‌کنند. توابع
\lr{read\_byte}
و
\lr{read\_byte\_1}
هشت نمونه را از کم‌ارزش‌ترین بیت در یک بایت جمع می‌کنند.

\subsection{توابع انتظار و بازنشانی}
توابع
\lr{B9\_wait\_for\_1}
و
\lr{C7\_wait\_for\_1}
گذرگاه را برای مدت ورودی آزاد می‌کنند. توابع
\lr{reset}
و
\lr{reset\_1}
پالس ۵۰۰ میکروثانیه‌ای پایین و ۵۰۰ میکروثانیه آزاد می‌سازند، اما حضور حسگر را نمونه‌برداری یا گزارش نمی‌کنند.

\subsection{توابع انتخاب شناسه}
تابع
\lr{get\_ID}
فرمان خواندن
\lr{ROM}
را روی گذرگاه اول اجرا و هشت بایت را در بافر داخلی قرار می‌دهد؛ فقط زمانی معتبر است که دقیقاً یک حسگر روی گذرگاه باشد. توابع
\lr{send\_ID}
و
\lr{send\_ID\_1}
پس از بازنشانی فرمان تطبیق
\lr{ROM}
و هشت بایت نشانی ورودی را می‌فرستند.

\subsection{توابع خواندن دما}
توابع
\lr{get\_temperature\_with\_ID}
و
\lr{get\_temperature\_with\_ID\_1}
حسگر را انتخاب، فرمان تبدیل را ارسال، مدت کوتاهی خط را آزاد، دوباره حسگر را انتخاب و نه بایت حافظه موقت را می‌خوانند. دو بایت نخست به عدد خام علامت‌دار تبدیل می‌شوند.
\lr{CRC}
حافظه موقت، بیت پایان تبدیل و حضور حسگر هنوز بررسی نمی‌شوند؛ بنابراین خروجی تابع تضمین صحت ندارد.

\section{درایور \lr{AM2315}}

\subsection{تابع \lr{am2315\_crc16}}
تابع داخلی،
\lr{CRC}
شانزده‌بیتی بازتاب‌یافته مودباس را با مقدار آغازین
\lr{0xFFFF}
و چندجمله‌ای
\lr{0xA001}
محاسبه می‌کند. ورودی داده و طول و خروجی
\lr{CRC}
است.

\subsection{تابع \lr{TEMP\_HUM}}
چهار ورودی آن دسته
\lr{I2C}،
بافر فرمان سه‌بایتی، بافر پاسخ هشت‌بایتی و آرایه دو عدد اعشاری خروجی است. ابتدا حسگر را با مرحله نشانی بیدار، سپس فرمان خواندن چهار رجیستر را ارسال و پاسخ را دریافت می‌کند. کد تابع، طول،
\lr{CRC}،
علامت دما و وضعیت
\lr{HAL}
بررسی می‌شوند. در موفقیت رطوبت و دما نوشته و رطوبت بازگردانده می‌شود؛ در شکست هر دو خروجی صفر و پرچم خطا یک می‌شود. این ماژول در مسیر فعال فعلی فراخوانی نمی‌شود.

\section{فیلتر \lr{MAF}}

\subsection{تابع \lr{MAF\_Init}}
اشاره‌گر یک نمونه فیلتر را می‌گیرد، معکوس طول پنجره را محاسبه، میانگین و همه ۴۰ خانه را صفر و اشاره‌گر نوشتن را روی خانه نخست قرار می‌دهد. باید پیش از به‌روزرسانی فراخوانی شود.

\subsection{تابع \lr{MAF\_Update}}
نمونه جدید را در خانه جاری می‌نویسد و اشاره‌گر را حلقوی جلو می‌برد. سپس همه ۴۰ خانه را جمع و میانگین را ذخیره و برمی‌گرداند. ورودی دوم داده جدید است. هزینه هر فراخوانی خطی است و نمونه‌های صفر اولیه در شروع روی نتیجه اثر دارند.

\subsection{تابع \lr{Average}}
اشاره‌گر آرایه و تعداد اعضا را می‌گیرد و میانگین حسابی برمی‌گرداند. اندازه صفر کنترل نشده و باعث تقسیم بر صفر می‌شود؛ فراخواننده موظف است اندازه مثبت بدهد. این تابع و کل ماژول در برنامه فعلی استفاده نمی‌شوند.
